\documentclass[../../main.tex]{subfiles}
\begin{document}
\section{交换向量场}\label{section:ch09-交换向量场}

设 $M$ 是光滑流形，$V,W\in\mathfrak{X}(M)$。我们说 $V$ 和 $W$ \textbf{交换}，若对每个光滑函数 $f$，$VWf=WVf$，等价地，若 $[V;W]\equiv 0$。若 $\theta$ 是光滑流，向量场 $W$ 称为在 $\theta$ 下\textbf{不变}，若对每个 $t$，$W$ 与自身是 $\theta_t$-相关的；更精确地说，这意味着对每个 $t$，$W|_{M_t}$ 与 $W|_{M_{-t}}$ 是 $\theta_t$-相关的，等价地，对 $\theta$ 定义域内的所有 $(t;p)$，$\mathrm{d}(\theta_t)_p(W_p)=W_{\theta_t(p)}$。下一个命题表明这两个概念密切相关。

\begin{theorem}\label{theorem:ch09-42}
对光滑流形 $M$ 上的光滑向量场 $V$ 和 $W$，下列条件等价：
\begin{enumerate}[(a)]
\item $V$ 和 $W$ 交换。
\item $W$ 在 $V$ 的流下不变。
\item $V$ 在 $W$ 的流下不变。
\end{enumerate}
\end{theorem}

\begin{proof}
假设 $V,W\in\mathfrak{X}(M)$，$\theta$ 表示 $V$ 的流。若 (b) 成立，则当 $(t;p)$ 在 $\theta$ 的定义域内时，$W_{\theta_t(p)}=\mathrm{d}(\theta_t)_p(W_p)$。将 $\mathrm{d}(\theta_{-t})_{\theta_{-t}(p)}$ 作用于两边，推出 $\mathrm{d}(\theta_{-t})_{\theta_{-t}(p)}(W_{\theta_{-t}(p)})=W_p$，由 Lie 导数的定义显然这直接推出 $[V;W]=L_VW=0$。同样的论证表明 (c) 推出 (a)。

为证明 (a) 推出 (b)，假设 $[V;W]=L_VW=0$。任取 $p\in M$，对 $t\in\mathscr{D}(p)$ 令 $X(t)=\mathrm{d}(\theta_{-t})_{\theta_{-t}(p)}(W_{\theta_{-t}(p)})$。\cref{proposition:ch09-41}表明 $X'(t)\equiv 0$。由于 $X(0)=W_p$，这推出对所有 $t\in\mathscr{D}(p)$，$X(t)=W_p$，将 $\mathrm{d}(\theta_t)_p$ 作用于两边即得到表明 $W$ 在 $\theta$ 下不变的恒等式。同样的证明也表明 (a) 推出 (c)。
\end{proof}

\begin{corollary}\label{corollary:ch09-43}
每个光滑向量场在其自身的流下不变。
\end{corollary}

\begin{proof}
将前一命题与 $[V;V]\equiv 0$ 这一事实结合使用。
\end{proof}

交换向量场最深刻的刻画在于它们各自的流之间的关系。下一个定理说两个向量场交换当且仅当它们的流交换。但在正式陈述定理之前，我们需要确切考察这意味着什么。假设 $V$ 和 $W$ 是 $M$ 上的光滑向量场，$\theta$ 和 $\psi$ 表示它们各自的流。若 $V$ 和 $W$ 都完备，说它们的流交换的意义是清楚的：就是对所有 $s,t\in\mathbb{R}$，$\theta_t\circ\psi_s=\psi_s\circ\theta_t$。然而，若 $V$ 或 $W$ 之一不完备，我们最多只能期望此等式对两边都有定义的所有 $s$ 和 $t$ 成立。不幸的是，即使向量场交换，它们的流也可能不在这种朴素意义下交换，因为存在交换向量场 $V$ 和 $W$ 以及 $t,s,p$ 的特殊选择，使得 $\theta_t\circ\psi_s(p)$ 和 $\psi_s\circ\theta_t(p)$ 都有定义但不相等（此类例子见习题 9-19）。问题在于：若 $\theta_t\circ\psi_s(p)$ 对 $t=t_0$ 和 $s=s_0$ 有定义，则由流域的性质，它必对包含 $0$ 和 $t_0$ 的某开区间中的所有 $t$ 有定义，但对 $s$ 的类似陈述不一定成立——可能存在 $0$ 和 $s_0$ 之间的某些 $s$ 值，使得从 $\psi_s(p)$ 出发的 $V$ 的积分曲线不能一直延拓到 $t=t_0$。

因此我们给出如下定义。若 $\theta$ 和 $\psi$ 是 $M$ 上的流，我们说 $\theta$ 和 $\psi$ \textbf{交换}，若对每个 $p\in M$ 下列条件成立：每当 $J$ 和 $K$ 是包含 $0$ 的开区间，使得表达式 $\theta_t\circ\psi_s(p)$ 或 $\psi_s\circ\theta_t(p)$ 之一对所有 $(s;t)\in J\times K$ 有定义时，两者都有定义且相等。对全局流，这等价于说对所有 $s$ 和 $t$，$\theta_t\circ\psi_s=\psi_s\circ\theta_t$。

\begin{theorem}\label{theorem:ch09-44}
光滑向量场交换当且仅当它们的流交换。
\end{theorem}

\begin{proof}
设 $V$ 和 $W$ 是光滑流形 $M$ 上的光滑向量场，$\theta$ 和 $\psi$ 表示它们各自的流。首先假设 $V$ 和 $W$ 交换。假设 $p\in M$，$J$ 和 $K$ 是包含 $0$ 的开区间，使得 $\psi_s\circ\theta_t(p)$ 对所有 $(s;t)\in J\times K$ 有定义。（将 $V$ 和 $W$ 互换的同样证明在另一表达式在此类矩形上有定义的假设下成立。）由\cref{theorem:ch09-42}，假设意味着 $V$ 在 $\psi$ 下不变。固定任意 $s\in J$，考虑由 $\gamma(t)=\psi_s\circ\theta_t(p)=\psi_s(\theta^{(p)}(t))$ 定义的曲线 $\gamma\colon K\to M$。此曲线满足 $\gamma(0)=\psi_s(p)$，其在 $t\in K$ 处的速度为
\begin{align*}
\gamma'(t)=\frac{\mathrm{d}}{\mathrm{d}t}\bigl(\psi_s(\theta^{(p)}(t))\bigr)=\mathrm{d}(\psi_s)(\theta^{(p)\prime}(t))=\mathrm{d}(\psi_s)(V_{\theta^{(p)}(t)})=V_{\gamma(t)}.
\end{align*}
因此，$\gamma$ 是从 $\psi_s(p)$ 出发的 $V$ 的积分曲线。故由唯一性，
\begin{align*}
\gamma(t)=\theta^{(\psi_s(p))}(t)=\theta_t(\psi_s(p)).
\end{align*}
这证明 $\theta$ 和 $\psi$ 交换。

反过来，假设流交换，令 $p\in M$。若选取 $\varepsilon>0$ 足够小，使得当 $|s|<\varepsilon$ 且 $|t|<\varepsilon$ 时 $\psi_s\circ\theta_t(p)$ 有定义，则假设保证对所有此类 $s$ 和 $t$，$\psi_s\circ\theta_t(p)=\theta_t\circ\psi_s(p)$。这可改写为
\begin{align*}
\psi^{(\theta_t(p))}(s)=\theta_t(\psi^{(p)}(s)).
\end{align*}
关于 $s$ 微分此关系，得
\begin{align*}
W_{\theta_t(p)}=\left.\frac{\mathrm{d}}{\mathrm{d}s}\right|_{s=0}\psi^{(\theta_t(p))}(s)=\left.\frac{\mathrm{d}}{\mathrm{d}s}\right|_{s=0}\theta_t(\psi^{(p)}(s))=\mathrm{d}(\theta_t)_p(W_p).
\end{align*}
将 $\mathrm{d}(\theta_{-t})_{\theta_{-t}(p)}$ 作用于两边，推出
\begin{align*}
\mathrm{d}(\theta_{-t})_{\theta_{-t}(p)}(W_{\theta_{-t}(p)})=W_p.
\end{align*}
关于 $t$ 微分并应用 Lie 导数的定义表明 $(L_VW)_p=0$。
\end{proof}

\subsection*{交换标架}

假设 $M$ 是光滑 $n$-流形。回忆 $M$ 的\textbf{局部标架}是 $n$ 元向量场 $(E_i)$，定义在开子集 $U\subseteq M$ 上，使得 $(E_i|_p)$ 在每个 $p\in U$ 处构成 $T_pM$ 的基。$M$ 的光滑局部标架 $(E_i)$ 称为\textbf{交换标架}，若对所有 $i$ 和 $j$，$[E_i;E_j]=0$。（交换标架被一些作者称为\textbf{完整标架}。）

\begin{example}\label{example:ch09-45}[交换与非交换标架]
\begin{enumerate}[(a)]
\item 交换标架最简单的例子是坐标标架。给定光滑流形 $M$ 的任意光滑坐标卡 $(U,(x^i))$，\eqref{align:ch08-10}表明坐标标架 $(\partial/\partial x^i)$ 是交换标架。

\item 由\eqref{align:ch08-03}定义的 $\mathbb{R}^2$ 上 $\mathbb{R}^2\smallsetminus\{0\}$ 的标架 $(E_1;E_2)$ 不是交换标架，因为直接计算表明
\begin{align*}
[E_1;E_2]=\frac{y}{r^2}\frac{\partial}{\partial x}-\frac{x}{r^2}\frac{\partial}{\partial y}\neq 0.
\end{align*}
\end{enumerate}
\end{example}

由于每个坐标标架都是交换标架，且 Lie 括号不变地定义，故光滑标架可在某光滑坐标卡中表示为坐标标架的必要条件是它为交换标架。因此，上面的计算表明 $(E_1;E_2)$ 不能关于任意光滑局部坐标选择表示为 $\mathbb{R}^2$ 的坐标标架。

下一个定理表明交换也是光滑标架局部可表示为坐标标架的充分条件。

\begin{theorem}\label{theorem:ch09-46}[交换向量场的标准型]
设 $M$ 是光滑 $n$-流形，$(V_1,\dots,V_k)$ 是开子集 $W\subseteq M$ 上线性无关的 $k$ 元光滑交换向量场。对每个 $p\in W$，存在以 $p$ 为中心的光滑坐标卡 $(U,(s^i))$，使得对 $i=1,\dots,k$，$V_i=\partial/\partial s^i$。若 $S\subseteq W$ 是嵌入的余维 $k$ 子流形，$p$ 是 $S$ 中使 $T_pS$ 与 $(V_1|_p,\dots,V_k|_p)$ 的张成互补的点，则坐标还可取得使 $S\cap U$ 是由 $s^1=\cdots=s^k=0$ 定义的切片。
\end{theorem}

\begin{proof}
任取 $p\in W$。若未给出子流形 $S$，就取任意光滑嵌入余维 $k$ 子流形 $S$，使其在 $p$ 处的切空间与 $(V_1|_p,\dots,V_k|_p)$ 的张成互补（例如，适当的坐标切片）。设 $(U,(x^i))$ 是以 $p$ 为中心的关于 $S$ 的切片坐标卡，$U\subseteq W$，$S\cap U$ 等于切片 $\{x\in U:x^1=\cdots=x^k=0\}$。我们的假设保证向量 $(V_1|_p,\dots,V_k|_p,\partial/\partial x^{k+1}|_p,\dots,\partial/\partial x^n|_p)$ 张成 $T_pM$。由于定理纯局部，不妨将 $V_1,\dots,V_k$ 视为 $U\subseteq\mathbb{R}^n$ 上的向量场，将 $S$ 视为 $U$ 中前 $k$ 个坐标为零的子集。此证明的基本思想与流出定理类似，只是需做一些额外工作以利用向量场交换的假设。

设 $\theta_i$ 表示 $V_i$ 的流（$i=1,\dots,k$）。存在 $\varepsilon>0$ 和 $p$ 在 $U$ 中的邻域 $Y$，使得复合 $(\theta_1)_{t_1}\circ(\theta_2)_{t_2}\circ\cdots\circ(\theta_k)_{t_k}$ 在 $Y$ 上有定义，且只要 $|t_1|,\dots,|t_k|$ 都小于 $\varepsilon$ 就将 $Y$ 映入 $U$。（为看到这一点，只需选 $\varepsilon_k>0$ 和 $U_k\subseteq U$ 使 $\theta_k$ 将 $(-\varepsilon_k,\varepsilon_k)\times U_k$ 映入 $U$，然后归纳选 $\varepsilon_i$ 和 $U_i$ 使 $\theta_i$ 将 $(-\varepsilon_i,\varepsilon_i)\times U_i$ 映入 $U_{i+1}$。取 $\varepsilon=\min\{\varepsilon_i\}$，$Y=U_1$ 即可。）

定义 $\Omega\subseteq\mathbb{R}^{n-k}$ 为
\begin{align*}
\Omega=\{(s^{k+1},\dots,s^n)\in\mathbb{R}^{n-k}:(0,\dots,0,s^{k+1},\dots,s^n)\in Y\},
\end{align*}
定义 $\Phi\colon(-\varepsilon,\varepsilon)^k\times\Omega\to U$ 为
\begin{align*}
\Phi(s^1,\dots,s^k,s^{k+1},\dots,s^n)=(\theta_1)_{s^1}\circ\cdots\circ(\theta_k)_{s^k}(0,\dots,0,s^{k+1},\dots,s^n).
\end{align*}
由构造，$\Phi(\{0\}\times\Omega)=S\cap Y$。

下面证明对 $i=1,\dots,k$，$\partial/\partial s^i$ 与 $V_i$ 是 $\Phi$-相关的。由于流 $\theta_i$ 交换，对任意 $i\in\{1,\dots,k\}$ 和任意 $s_0\in(-\varepsilon,\varepsilon)^k\times\Omega$，有
\begin{align*}
\mathrm{d}\Phi_{s_0}\left(\frac{\partial}{\partial s^i}\bigg|_{s_0}\right)f
&=\left.\frac{\partial}{\partial s^i}\right|_{s_0}f\bigl(\Phi(s^1,\dots,s^n)\bigr)\\
&=\left.\frac{\partial}{\partial s^i}\right|_{s_0}f\Bigl((\theta_1)_{s^1}\circ\cdots\circ(\theta_k)_{s^k}(0,\dots,0,s^{k+1},\dots,s^n)\Bigr)\\
&=\left.\frac{\partial}{\partial s^i}\right|_{s_0}f\Bigl((\theta_i)_{s^i}\circ(\theta_1)_{s^1}\circ\cdots\circ(\theta_{i-1})_{s^{i-1}}\circ(\theta_{i+1})_{s^{i+1}}\circ\cdots\circ(\theta_k)_{s^k}(0,\dots,0,s^{k+1},\dots,s^n)\Bigr).
\end{align*}
对任意 $q\in M$，$t\mapsto(\theta_i)_t(q)$ 是 $V_i$ 的积分曲线，故最后一个表达式等于 $V_i|_{\Phi(s_0)}f$，证明断言。

接下来证明 $\mathrm{d}\Phi_0$ 可逆。上面的计算表明
\begin{align*}
\mathrm{d}\Phi_0\left(\frac{\partial}{\partial s^i}\bigg|_0\right)=V_i|_p,\qquad i=1,\dots,k.
\end{align*}
另一方面，由于 $\Phi(0,\dots,0,s^{k+1},\dots,s^n)=(0,\dots,0,s^{k+1},\dots,s^n)$，直接推出
\begin{align*}
\mathrm{d}\Phi_0\left(\frac{\partial}{\partial s^i}\bigg|_0\right)=\frac{\partial}{\partial x^i}\bigg|_p,\qquad i=k+1,\dots,n.
\end{align*}
由此 $\mathrm{d}\Phi_0$ 将 $T_0\mathbb{R}^n$ 的基 $(\partial/\partial s^1|_0,\dots,\partial/\partial s^n|_0)$ 映为 $T_pM$ 的基 $(V_1|_p,\dots,V_k|_p,\partial/\partial x^{k+1}|_p,\dots,\partial/\partial x^n|_p)$。由反函数定理，$\Phi$ 在 $0$ 邻域上是微分同胚，$\varphi=\Phi^{-1}$ 是光滑坐标映射，它将 $\partial/\partial s^i$ 映为 $V_i$（$i=1,\dots,k$），将 $S$ 映为切片 $s^1=\cdots=s^k=0$。
\end{proof}

正如单向量场的情形，\cref{theorem:ch09-46}的证明提示了一种寻找显式坐标使一组交换向量场化为标准型的技巧，只要它们的流可显式求出。该方法可总结如下：从一个 $(n-k)$ 维子流形 $S$ 开始，其在 $p$ 处的切空间与 $(V_1|_p,\dots,V_k|_p)$ 的张成互补。然后通过从 $S$ 中任意点出发、依次跟随 $k$ 个流任意时间来定义 $\Phi$。由于流交换，施加顺序无关。一个例子有助于澄清该过程。

\begin{example}\label{example:ch09-47}
考虑 $\mathbb{R}^2$ 上如下两个向量场：
\begin{align*}
V=x\frac{\partial}{\partial y}-y\frac{\partial}{\partial x},\qquad W=x\frac{\partial}{\partial x}+y\frac{\partial}{\partial y}.
\end{align*}
计算表明 $[V;W]=0$。\cref{example:ch09-08}表明 $V$ 的流为 $\theta_t(x,y)=(x\cos t-y\sin t,x\sin t+y\cos t)$，容易验证 $W$ 的流为 $\psi_t(x,y)=(e^t x,e^t y)$。在 $p=(1;0)$ 处，$V_p$ 和 $W_p$ 线性无关。由于此情形下 $k=n=2$，可取子集 $S$ 为单点 $\{(1;0)\}$，定义 $\Phi\colon\mathbb{R}^2\to\mathbb{R}^2$ 为
\begin{align*}
\Phi(s;t)=\psi_t\circ\theta_s(1;0)=(e^t\cos s,e^t\sin s).
\end{align*}
在此情形下，可在 $(1;0)$ 邻域内显式解出 $(s;t)=\Phi^{-1}(x;y)$，得到坐标映射
\begin{align*}
(s;t)=\bigl(\tan^{-1}(y/x),\log\sqrt{x^2+y^2}\bigr).
\end{align*}
\end{example}

\end{document}